\chapter{Bridges, Live Link and file formats}
\label{ch:bridges}

\section{Live Link}
\begin{figure}[H]
  \centering
  \resizebox{0.95\textwidth}{!}{% The Live Link sequence: discovery, connection, the sync, then changes as they happen.
\begin{tikzpicture}[node distance=4mm]
  \node[box, fill=blAmber!15, draw=blAmber, minimum width=34mm] (b) {Blendity\\ \footnotesize \code{src/editor/livelink.*}};
  \node[box, right=100mm of b, minimum width=34mm] (c) {A bridge\\ \footnotesize Blender / Unity / Godot / Unreal};
  \coordinate (bb) at ($(b.south)+(0,-96mm)$);
  \coordinate (cb) at ($(c.south)+(0,-96mm)$);
  \draw[blGrey, line width=0.5pt] (b.south) -- (bb);
  \draw[blGrey, line width=0.5pt] (c.south) -- (cb);
  \foreach \y/\dir/\txt in {
    8/r/{writes \code{\textasciitilde/.blendity/livelink.json}: the port, pid, scene and inbox},
    19/l/{reads the file every 2 s while not connected; connects to the port (TCP, 127.0.0.1 only)},
    30/l/{\code{hello}: the app, the bridge's version, its caps (\code{skins})},
    41/r/{\code{sync} \ldots\ \code{sync\_end}: every object, mesh (a glTF binary), material, light, camera, the world},
    54/r/{each frame, at most 60 a second: the difference since the last comparison - \code{transform}, \code{pose}, \code{deform}, \code{mesh}, \code{material}, \code{texture}, \code{clip}, \code{delete}, \code{play} \ldots},
    68/r/{a slow client's queue is coalesced by key (a newer \code{transform} replaces the queued one), never past a \code{mesh} or a \code{delete}},
    80/l/{Send to Blendity (Blender, Unreal): a snapshot folder into the inbox, imported within a second},
    91/r/{\code{livelink stop} or exit: the discovery file is removed}}{
    \ifx\dir r
      \draw[flow] ($(b.south)+(0,-\y mm)$) -- node[lbl, above, text width=92mm, align=center] {\txt} ($(c.south)+(0,-\y mm)$);
    \else
      \draw[flow] ($(c.south)+(0,-\y mm)$) -- node[lbl, above, text width=92mm, align=center] {\txt} ($(b.south)+(0,-\y mm)$);
    \fi
  }
\end{tikzpicture}
}
  \caption{The Live Link sequence (\code{docs/LIVELINK.md}, \adr{0055}, \adr{0086}). One way out: nothing a client sends is read but its hello.}
  \label{fig:livelink}
\end{figure}

Live Link sends the scene open in Blendity to other programs on the same computer as it changes. The transport is a
TCP server on 127.0.0.1 (\code{src/core/net\_server.*}, no library); the framing is one JSON object a line with a
binary attachment of N bytes when the line ends with \code{"bin":N}; meshes go as glTF 2.0 binaries
\cite{Khronos2021Gltf} from Blendity's own exporter in glTF's axes (right-handed, Y up, metres: X mirrored), materials as
glTF material objects with the toon fields as an extension, textures by their pixels, skins once at rest and then
poses and weights, clips baked as the engines' own, colliders and LOD levels, the world and the camera's look, VRChat's
components where the SDK is there. What changed is found by comparing the scene with what was last sent, once a frame
while a client is connected, so a drag streams while it happens and undo, Stop and a reload need no code of their
own; each client has its own queue and sender thread, coalesced by key so a slow client skips to the present. With four
engines connected, a 100,000-vertex mesh changing and a 50-bone rig posing every frame, it takes a millisecond or two
of a frame (\rnd{382}).

\section{The bridges}
\code{bridges/} holds one package per engine, Blendity's own code under the same GPL: a Blender add-on (Python), a
Unity package (\code{com.blendity.livelink}, C\#, with VRChat stand-ins and SDK checks), a Godot add-on (GDScript) and
an Unreal plugin (C++), each with a smoke test, an auto-connect test and a live session test, packaged by
\code{tools/package\_bridges.py} and installed from Window $>$ Live Link $>$ Install a Bridge. The Blender and Unreal
bridges also send to Blendity: a snapshot folder into the inbox the discovery file names, imported within a second.
\code{tools/livelink\_fake\_server.py} serves a recorded session for a bridge's tests, and \code{livelink\_dump.py} is
the reading loop in fifty lines of Python.

\section{File formats}
Export shares one preparation step and a writer per format (\code{src/scene/export.cpp}): OBJ + MTL, FBX 7.4 as a
record tree with Blender's binary encoding \cite{FBXBlender2026} and an ASCII one, glTF 2.0 (\code{.glb} or
\code{.gltf}, embedded or copied textures, morph targets, skins and sampled channels), STL, PLY, USD text
\cite{USD2026}; options as Blender's (Selected Only, Apply Modifiers, Scale, Up Axis, Triangulate, Keep Transforms,
Armatures and Animation, Colliders and LOD levels named for Unity, Unreal or Godot). Every format was checked both ways
against Blender 5.2 (R-06). Import reads the same formats plus MMD's PMX / PMD / VMD \cite{MMD2026}, VRM
\cite{MToon2026}, BVH, Collada, and \code{.blend} through the installed Blender, with an Import Preview that shows the
model under a chosen shading before anything is written (\adr{0047}).

\section{Unity's project files}
Rounds 369--408 made Blendity read and write Unity's own files: \code{.unitypackage} unpacked into the project with
GUIDs kept; prefabs and scenes brought in with nested prefabs and overrides, every other component as an Engine
Component editable under Unity's names (\adr{0101}); \code{.mat} files converted by shader (Standard, URP, lilToon,
Poiyomi, MToon, UTS2, Filamented); a \code{.unity} opened as the scene with the look Unity gives it (fog modes, skybox,
camera flags, post-processing, lights' shadows, reflection probes, \adr{0099}; Unity's own lightmaps and probes read
by the installed Unity in the background, \adr{0098}); VRChat's components as Blendity's (\adr{0100}); Unity's
\code{.anim} clips read and written, the Unity file never written over (\adr{0102}); humanoid retargeting as Mecanim
reads a T and an A pose (\adr{0103}); Export Unity Package writing the work back under the same GUIDs. The parity is
measured: \code{tools/unity\_compare.py} renders the same stations in Unity (batch mode) and in Blendity and scores them
(MAE, PSNR, SSIM, a luminance grid's correlation), and \code{tools/unity\_e2e.py} runs the whole fresh-project flow;
an avatar passes at every station, a world is near the target (\code{docs/hci/compare-unity-vrchat-2026-11.md}).

\section{Godot, Blender, Unreal}
Godot scenes (\code{.tscn}) come in with their nodes, primitives, materials, lights, cameras and collision shapes and go
back with the objects' transforms in their own lines, scripts kept (\rnd{373}, \cite{GodotDocs2026}). Blender files come
in through the installed Blender's glTF exporter with lights, cameras, actions and custom properties (\rnd{374}). Unreal's
bridge sends a level through its FBX exporter with each actor's place and Blueprint class (\cite{UnrealDocs2026}).

\section{MoonRay}
DreamWorks' MoonRay \cite{MoonRay2026} is an optional external renderer (\adr{0041}): built by the \code{moonray.yml}
workflow on Linux, run as a program of its own (in WSL on Windows), fed a scene written by Blendity, its picture read
back into the Render window. It is never linked and never shipped with Blendity.
